\section{Additional tables}\label{app:tables}
Unless a caption says otherwise, every table in the appendices uses v3 with its walk-forward hyper-parameters, the full run on the frozen snapshot and the EDGAR universe (cutoffs 2014--2024, eight-quarter horizon). The online supplementary material also reports results from the earlier evaluation's design: the four-quarter forecasts use S\&P~500 members (today's, a sample selected on survival, and those of 2020), and the 2020 analyses use firms alive on 20 February 2020. Their samples, and therefore some AUCs and counts, differ from those of the main text.

\begin{table}[!htbp]
\centering\scriptsize
\caption{Calibrated parameters}
\label{tab:params}
\begin{tabular}{lp{0.5\linewidth}l}
\toprule
Parameter & Estimate & Window \\
\midrule
Seasonal factors $S_q$ & median ratio to centred 4-quarter moving average & last 32 quarters \\
Trend growth $\mu$ & median log growth, shrunk 50\% toward 1.2\%/quarter & last 12 quarters \\
AR coefficient $\phi$, volatility $\sigma_g$ & OLS; median absolute deviation of residuals & last 32 quarters \\
Tail index & Student-$t$ maximum likelihood on standardized residuals & growth residuals \\
Cyclicality $\beta$ & regression of revenue on real-time nominal GDP growth, shrunk 50\% toward 1.5 & last 40 quarters \\
Recession entry and exit & NBER recessions since 1947 & FRED \\
Fixed costs $F$, variable ratio $v$ & v3: posterior of a regression of cash costs on revenue (four-quarter changes), prior s.d.\ $\tau$ walk-forward; v2: penalized regression, prior 35\% (v2-pre) or 5\% (v2-rev) fixed & last 16 quarters \\
Fixed-cost adjustment $\lambda$ & v3: partial adjustment to permanent revenue changes, walk-forward & --- \\
Margin shock $m_t$ & AR(1) of cost-regression residuals; correlation with demand & last 16 quarters \\
DSO, DIO, DPO, deposit days & averages of balance to revenue or costs & last 4 quarters \\
Depreciation, capex, SBC ratios & medians & last 8 quarters \\
Coupon & interest expense over debt & last 4 quarters \\
Rollover share & two-year debt trend relative to maturities & last 8 quarters \\
Discretionary financing volatility & robust s.d.\ of quarterly debt changes & last 20 quarters \\
Maturity ladder & XBRL maturity schedule, else six-year average life & latest 10-K \\
Credit facility & XBRL facility size, else 8\% of revenue (assumption) & latest filing \\
Minimum cash & half the 10th-percentile cash balance & last 12 quarters \\
Short-rate dynamics & Vasicek fit to the 3-month Treasury bill & since 2000 \\
\bottomrule
\end{tabular}
\end{table}

\begin{table}[!htbp]
\centering\scriptsize
\caption{Forecast encompassing (v3)}
\label{tab:encompassing}
\resizebox{\linewidth}{!}{%
\begin{tabular}{llrrrrrr}
\toprule
Cash measure & Specification & 4q margin & Twin & 16q margin & Shrunk margin & Regression + timing & $R^2$ \\
\midrule
Operating cash flow & Departures: twin only &  & 0.43 (13.9)$^{***}$ &  &  &  & 0.21 \\
 & Departures: simple rules only &  &  & 0.42 (11.2)$^{***}$ & 0.38 (7.0)$^{***}$ & 0.34 (11.2)$^{***}$ & 0.32 \\
 & Departures: all &  & 0.23 (8.5)$^{***}$ & 0.33 (8.8)$^{***}$ & 0.31 (5.6)$^{***}$ & 0.27 (9.7)$^{***}$ & 0.36 \\
 & Levels: simple rules only & 0.15 (2.4)$^{**}$ &  & 0.34 (8.9)$^{***}$ & 0.05 (0.8) & 0.31 (10.7)$^{***}$ & 0.67 \\
 & Levels: all & 0.10 (1.7)$^{*}$ & 0.21 (7.6)$^{***}$ & 0.28 (7.1)$^{***}$ & 0.03 (0.5) & 0.25 (9.3)$^{***}$ & 0.69 \\
\addlinespace
Free cash flow & Departures: twin only &  & 0.45 (15.4)$^{***}$ &  &  &  & 0.23 \\
 & Departures: simple rules only &  &  & 0.40 (10.6)$^{***}$ & 0.61 (10.3)$^{***}$ & 0.29 (11.8)$^{***}$ & 0.31 \\
 & Departures: all &  & 0.24 (9.4)$^{***}$ & 0.31 (8.5)$^{***}$ & 0.47 (8.5)$^{***}$ & 0.23 (9.6)$^{***}$ & 0.36 \\
 & Levels: simple rules only & 0.06 (0.8) &  & 0.33 (8.7)$^{***}$ & 0.20 (2.7)$^{***}$ & 0.26 (10.3)$^{***}$ & 0.61 \\
 & Levels: all & 0.02 (0.3) & 0.22 (8.4)$^{***}$ & 0.27 (7.2)$^{***}$ & 0.16 (2.4)$^{**}$ & 0.20 (8.6)$^{***}$ & 0.63 \\
\bottomrule
\end{tabular}}

\par\smallskip
\begin{minipage}{\linewidth}\footnotesize
Departures: OLS of (realized $-$ four-quarter margin) on each model's departure from the same rule. Levels: OLS of realized cash flow on the forecasts themselves \citep{fair1990}. All variables winsorized at the 1st and 99th percentiles; year effects; standard errors clustered by firm; $t$-statistics in parentheses. $^{*}$, $^{**}$, $^{***}$: $p<0.10, 0.05, 0.01$.
\end{minipage}
\end{table}

\begin{table}[!htbp]
\centering\footnotesize
\caption{Out-of-sample combinations, implementation of the earlier evaluation (v3)}
\label{tab:hybrid}
\resizebox{\linewidth}{!}{%
\begin{tabular}{lrrrr}
\toprule
 & \multicolumn{2}{c}{Operating cash flow} & \multicolumn{2}{c}{Free cash flow} \\
\cmidrule(lr){2-3}\cmidrule(lr){4-5}
Out-of-sample prediction & Mean error & Median error & Mean error & Median error \\
\midrule
Margin, 4 quarters & 13.97 & 8.15 & 15.08 & 8.69 \\
Margin, 16 quarters & 12.84 & 7.14 & 13.97 & 7.56 \\
Regression + revenue change & 15.70 & 9.43 & 17.22 & 10.10 \\
Twin alone & 14.80 & 8.41 & 15.84 & 8.83 \\
4-quarter margin + $w$(twin $-$ rule) & 13.08 & 7.44 & 13.94 & 7.74 \\
16-quarter margin + $w$(4-quarter margin $-$ rule) & 12.65 & 7.17 & 13.66 & 7.56 \\
16-quarter margin + $w$(regression $-$ rule) & 12.49 & 7.14 & 13.64 & 7.50 \\
16-quarter margin + $w$(regression with timing $-$ rule) & 12.52 & 7.19 & 13.70 & 7.53 \\
16-quarter margin + $w$(twin $-$ rule) & 12.37 & 6.97 & 13.44 & 7.39 \\
\midrule
\multicolumn{5}{l}{\emph{Difference in mean error (negative: the first is better), firm-clustered CI}} \\
16q margin + twin vs.\ 16q margin & \multicolumn{2}{r}{$-$0.46 [$-$0.63, $-$0.29]$^{***}$} & \multicolumn{2}{r}{$-$0.52 [$-$0.72, $-$0.33]$^{***}$} \\
16q margin + twin vs.\ 16q margin + 4q margin & \multicolumn{2}{r}{$-$0.27 [$-$0.46, $-$0.05]$^{**}$} & \multicolumn{2}{r}{$-$0.22 [$-$0.41, $-$0.02]$^{**}$} \\
16q margin + twin vs.\ 16q margin + regression & \multicolumn{2}{r}{$-$0.11 [$-$0.25, 0.06]} & \multicolumn{2}{r}{$-$0.20 [$-$0.37, $-$0.02]$^{**}$} \\
16q margin + twin vs.\ 16q margin + regression with timing & \multicolumn{2}{r}{$-$0.14 [$-$0.28, 0.03]$^{*}$} & \multicolumn{2}{r}{$-$0.26 [$-$0.43, $-$0.07]$^{***}$} \\
4q margin + twin vs.\ 4q margin & \multicolumn{2}{r}{$-$0.88 [$-$1.02, $-$0.72]$^{***}$} & \multicolumn{2}{r}{$-$1.14 [$-$1.34, $-$0.96]$^{***}$} \\
Twin vs.\ 16q margin & \multicolumn{2}{r}{+1.97 [1.48, 2.45]$^{***}$} & \multicolumn{2}{r}{+1.88 [1.30, 2.41]$^{***}$} \\
\bottomrule
\end{tabular}}

\par\smallskip
\begin{minipage}{\linewidth}\footnotesize
Each combination's weight ($w$ in the table) is chosen on earlier years. Errors in percentage points of trailing revenue (mean: winsorized at full-sample bounds). Differences: firm-clustered bootstrap CI.
\end{minipage}
\end{table}

\begin{table}[!htbp]
\centering\scriptsize
\caption{Transmission by firm characteristics and by revenue shock}
\label{tab:heterogeneity}
\begin{tabular}{llrrrrrr}
\toprule
 & & & \multicolumn{3}{c}{Incremental $\rho$} & \multicolumn{2}{c}{Median error} \\
\cmidrule(lr){4-6}\cmidrule(lr){7-8}
Split by & Group & Firm-years & Twin vs.\ 4q & 16q vs.\ 4q & Twin vs.\ regr. & 4q margin & Twin \\
\midrule
Net debt / EBITDA & low & 5{,}311 & 0.42 & 0.43 & 0.53 & 7.4 & 7.8 \\
 & mid & 5{,}311 & 0.42 & 0.46 & 0.56 & 5.8 & 6.4 \\
 & high & 5{,}311 & 0.35 & 0.34 & 0.44 & 9.9 & 11.4 \\
\addlinespace
Fixed-cost share & low & 5{,}311 & 0.40 & 0.43 & 0.52 & 6.8 & 7.0 \\
 & mid & 5{,}311 & 0.40 & 0.42 & 0.51 & 7.0 & 7.4 \\
 & high & 5{,}311 & 0.37 & 0.36 & 0.46 & 8.8 & 10.7 \\
\addlinespace
Debt due / cash & low & 5{,}311 & 0.37 & 0.37 & 0.49 & 8.9 & 9.7 \\
 & mid & 5{,}311 & 0.38 & 0.39 & 0.48 & 7.1 & 7.8 \\
 & high & 5{,}311 & 0.41 & 0.46 & 0.53 & 6.6 & 7.3 \\
\addlinespace
Revenue over the window & fell $>$10\% & 2{,}716 & 0.27 & 0.31 & 0.51 & 7.6 & 9.4 \\
 & fell 0-10\% & 2{,}756 & 0.37 & 0.39 & 0.30 & 6.1 & 6.7 \\
 & rose 0-10\% & 3{,}745 & 0.40 & 0.46 & 0.44 & 5.3 & 6.2 \\
 & rose $>$10\% & 6{,}716 & 0.43 & 0.42 & 0.55 & 10.0 & 10.1 \\
\bottomrule
\end{tabular}

\par\smallskip
\begin{minipage}{\linewidth}\footnotesize
Terciles of each characteristic at the cutoff; revenue over the window relative to trailing revenue. Incremental $\rho$ as in Table~\ref{tab:transmission}. The simulator's incremental correlation against the timing regression is \HetLevLowTwinReg{} in the lowest leverage tercile and \HetLevHighTwinReg{} in the highest; among the \HetDownN{} firm-years with revenue more than 10\% below the trailing level, its correlation against the four-quarter margin (\HetDownTwin) is below the 16-quarter margin's (\HetDownSixteen).
\end{minipage}
\end{table}

\begin{table}[!htbp]
\centering\scriptsize
\caption{Bankruptcy petitions (8-K Item 1.03), all non-financial US-GAAP filers}
\label{tab:bankruptcy}
\resizebox{\linewidth}{!}{%
\begin{tabular}{lllll}
\toprule
 & \multicolumn{2}{c}{Panel 2014--2024, within 8 quarters} & \multicolumn{2}{c}{20 Feb 2020, to end-2021} \\
\cmidrule(lr){2-3}\cmidrule(lr){4-5}
 & All (420 events) & Already weak (335) & All (52) & Already weak (37) \\
\midrule
Twin: health score (panel: baseline; 2020: stress) & 0.87 [0.85, 0.89] & 0.79 [0.76, 0.82] & 0.87 [0.82, 0.91] & 0.77 [0.70, 0.83] \\
Twin: breach probability & 0.84 [0.81, 0.87] & 0.78 [0.74, 0.81] & 0.74 [0.67, 0.81] & 0.64 [0.56, 0.73] \\
Health formula, no simulation & 0.87 [0.85, 0.89] & 0.77 [0.74, 0.80] & 0.87 [0.82, 0.91] & 0.76 [0.70, 0.82] \\
\addlinespace
Altman Z$''$ & 0.82 [0.79, 0.84] & 0.66 [0.63, 0.70] & 0.81 [0.76, 0.86] & 0.64 [0.58, 0.71] \\
Ohlson O-score & 0.86 [0.83, 0.88] & 0.77 [0.74, 0.79] & 0.79 [0.72, 0.87] & 0.73 [0.64, 0.83] \\
Zmijewski score & 0.82 [0.80, 0.84] & 0.75 [0.72, 0.78] & 0.78 [0.71, 0.85] & 0.75 [0.67, 0.83] \\
Interest coverage & 0.80 [0.78, 0.82] & 0.64 [0.61, 0.67] & 0.78 [0.72, 0.83] & 0.62 [0.54, 0.69] \\
Net debt / EBITDA & 0.80 [0.77, 0.82] & 0.67 [0.64, 0.69] & 0.77 [0.70, 0.82] & 0.63 [0.54, 0.71] \\
Cash / revenue & 0.59 [0.56, 0.63] & 0.57 [0.53, 0.61] & 0.63 [0.55, 0.70] & 0.63 [0.54, 0.71] \\
Current ratio & 0.68 [0.65, 0.71] & 0.61 [0.57, 0.64] & 0.72 [0.65, 0.78] & 0.63 [0.54, 0.71] \\
Static stress test & 0.75 [0.72, 0.78] & 0.70 [0.66, 0.73] & 0.75 [0.67, 0.82] & 0.68 [0.59, 0.76] \\
Revenue volatility & 0.67 [0.62, 0.70] & 0.59 [0.55, 0.63] & 0.63 [0.54, 0.71] & 0.56 [0.45, 0.65] \\
Fixed-cost share & 0.63 [0.59, 0.67] & 0.61 [0.57, 0.64] & 0.70 [0.63, 0.77] & 0.65 [0.56, 0.73] \\
Size (log revenue) & 0.63 [0.60, 0.66] & 0.52 [0.48, 0.56] & 0.57 [0.51, 0.64] & 0.42 [0.34, 0.51] \\
\addlinespace
Trailing FCF margin & 0.78 [0.74, 0.81] & 0.58 [0.53, 0.62] & 0.77 [0.70, 0.84] & 0.52 [0.41, 0.62] \\
Capex / revenue & 0.51 [0.47, 0.56] & 0.48 [0.44, 0.53] & 0.49 [0.38, 0.61] & 0.49 [0.38, 0.60] \\
\bottomrule
\end{tabular}}

\par\smallskip
\begin{minipage}{\linewidth}\footnotesize
AUC with firm-clustered 95\% CI. Panel: baseline health score at each cutoff; 2020: stress health score on 20 February 2020. Already weak: bottom quintile of Altman's Z$''$ in the cutoff year or negative trailing FCF margin.
\end{minipage}
\end{table}

\begin{table}[!htbp]
\centering\scriptsize
\caption{Distress events over the next eight quarters, 2014--2024}
\label{tab:panel}
\resizebox{\linewidth}{!}{%
\begin{tabular}{lrrrrlrrl}
\toprule
 & & & \multicolumn{3}{c}{Univariate AUC} & \multicolumn{3}{c}{Out-of-sample logit AUC} \\
\cmidrule(lr){4-6}\cmidrule(lr){7-9}
Event (next 8 quarters) & Firm-years & Events & Twin & Formula & Best benchmark & Benchmarks & + Twin & $\Delta$ [95\% CI] \\
\midrule
Severe cash burn & 15{,}933 & 2{,}368 & 0.76 & 0.74 & 0.82 (trailing FCF margin) & 0.87 & 0.87 & +0.004 [0.002, 0.005]$^{***}$ \\
Severe abnormal burn & 14{,}152 & 3{,}241 & 0.54 & 0.53 & 0.62 (revenue volatility) & 0.69 & 0.68 & $-$0.006 [$-$0.010, $-$0.003]$^{***}$ \\
Coverage collapse & 8{,}790 & 1{,}790 & 0.60 & 0.60 & 0.66 (interest coverage) & 0.69 & 0.70 & +0.004 [0.001, 0.007]$^{***}$ \\
Dividend cut & 8{,}686 & 716 & 0.70 & 0.69 & 0.70 (Ohlson O-score) & 0.76 & 0.76 & +0.007 [0.002, 0.011]$^{***}$ \\
Bankruptcy petition & 21{,}084 & 420 & 0.87 & 0.87 & 0.86 (Ohlson O-score) & 0.87 & 0.88 & +0.011 [0.005, 0.018]$^{***}$ \\
\bottomrule
\end{tabular}}

\par\smallskip
\begin{minipage}{\linewidth}\footnotesize
Twin: baseline health score. Out-of-sample logits trained for each test year on cutoffs at least three years earlier; ``Benchmarks'' contains all thirteen benchmarks and the formula, ``+ Twin'' adds the simulator's baseline health, breach probability and recession-scenario health. $\Delta$: firm-clustered CI.
\end{minipage}
\end{table}

\begin{table}[!htbp]
\centering\scriptsize
\caption{Liquidity events given the realized revenue}
\label{tab:decisions}
\resizebox{\linewidth}{!}{%
\begin{tabular}{lrrrrrrrrl}
\toprule
 & & & \multicolumn{4}{c}{Univariate AUC} & \multicolumn{3}{c}{Out-of-sample logit AUC} \\
\cmidrule(lr){4-7}\cmidrule(lr){8-10}
Event in the next 8 quarters & Firm-years & Rate & Twin & Static runway & Formula & FCF margin & Rule + scores & + Twin & $\Delta$ [95\% CI] \\
\midrule
Cash falls below half & 15{,}866 & 29\% & 0.55 & 0.52 & 0.56 & 0.54 & 0.56 & 0.56 & +0.009 [0.005, 0.013]$^{***}$ \\
New borrowing $>$ 5\% of revenue & 14{,}780 & 47\% & 0.58 & 0.56 & 0.52 & 0.49 & 0.58 & 0.61 & +0.022 [0.013, 0.030]$^{***}$ \\
Dividend cut by half & 7{,}714 & 8\% & 0.66 & 0.59 & 0.69 & 0.62 & 0.68 & 0.69 & +0.008 [$-$0.000, 0.015]$^{*}$ \\
Capex cut by a quarter & 15{,}362 & 20\% & 0.64 & 0.65 & 0.57 & 0.60 & 0.64 & 0.67 & +0.033 [0.025, 0.041]$^{***}$ \\
\bottomrule
\end{tabular}}

\par\smallskip
\begin{minipage}{\linewidth}\footnotesize
Twin: share of simulated paths with the event. Static runway: the trailing FCF margin on realized revenue, minus payouts and unrefinanced maturities (same inputs). Out-of-sample: ridge logits on within-year percentile ranks, trained on cutoffs up to $Y-3$; ``Rule + scores'' combines the static runway, the health formula and the FCF margin. $\Delta$: firm-clustered bootstrap CI.
\end{minipage}
\end{table}

\begin{table}[!htbp]
\centering\footnotesize
\begin{threeparttable}
\caption{Feedback loops evaluated on the events they are meant to produce}
\label{tab:loops}
\begin{tabular}{lrrrrrr}
\toprule
Event & Full twin & No R1 & No R2 & No B1 & No B2 & No loops \\
\midrule
Cash falls below half & 0.55 & 0.55 & 0.50$^{***}$ & 0.55$^{***}$ & 0.55 & 0.50$^{***}$ \\
New borrowing $>$ 5\% of revenue & 0.58 & 0.58 & 0.58$^{***}$ & 0.59$^{***}$ & 0.58$^{***}$ & 0.58 \\
Dividend cut by half & 0.66 & 0.66$^{*}$ & 0.66 & 0.50$^{***}$ & 0.66 & 0.50$^{***}$ \\
Capex cut by a quarter & 0.64 & 0.64$^{**}$ & 0.64 & 0.61$^{***}$ & 0.64 & 0.61$^{***}$ \\
\bottomrule
\end{tabular}

\begin{tablenotes}\footnotesize
\item AUC of the event probability of the full simulator and of the simulator with one loop switched off. Stars: paired difference with the full simulator.
\end{tablenotes}
\end{threeparttable}
\end{table}

\begin{table}[!htbp]
\centering\scriptsize
\caption{Liquidity events at other thresholds (given realized revenue)}
\label{tab:thresholds}
\begin{tabular}{lrrrr}
\toprule
Event & Threshold & Rate & Twin AUC & No loops AUC \\
\midrule
Cash falls below share of opening & 30\% & 14\% & 0.55 & 0.50 \\
Cash falls below share of opening & 50\% & 29\% & 0.55 & 0.50 \\
Cash falls below share of opening & 70\% & 49\% & 0.58 & 0.59 \\
New borrowing above share of revenue & 3\% & 53\% & 0.59 & 0.59 \\
New borrowing above share of revenue & 5\% & 47\% & 0.58 & 0.58 \\
New borrowing above share of revenue & 10\% & 35\% & 0.57 & 0.58 \\
Dividend cut by at least & 25\% & 12\% & 0.66 & 0.50 \\
Dividend cut by at least & 50\% & 8\% & 0.66 & 0.50 \\
Capex cut by at least & 25\% & 20\% & 0.64 & 0.61 \\
Capex cut by at least & 50\% & 7\% & 0.62 & 0.59 \\
\bottomrule
\end{tabular}

\end{table}

\begin{table}[!htbp]
\centering\scriptsize
\caption{Holm-adjusted $p$-values of the secondary comparisons (winsorization limits from all years; see the limitations)}
\label{tab:holm}
\begin{tabular}{lr}
\toprule
Secondary comparison & Holm-adjusted $p$ \\
\midrule
Liquidity: capex cut, twin added & $<$0.001 \\
Liquidity: cash halves, twin added & $<$0.001 \\
Liquidity: new borrowing, twin added & $<$0.001 \\
Combination: OCF twin vs. 16q + 4q & $<$0.001 \\
Combination: FCF twin vs. 16q + 4q & $<$0.001 \\
Combination: FCF twin vs. 16q + regression & $<$0.001 \\
Combination: FCF twin vs. 16q + regression with timing & $<$0.001 \\
Distress: abnormal burn, twin beyond formula & $<$0.001 \\
Distress: severe burn, twin beyond formula & $<$0.001 \\
Distress: bankruptcy, twin beyond formula & 0.032 \\
Distress: coverage collapse, twin beyond formula & 0.032 \\
Distress: dividend cut, twin beyond formula & 0.032 \\
Bankruptcy AUC: twin minus formula, weak firms & 0.060 \\
Combination: OCF twin vs. 16q + regression with timing & 0.080 \\
Combination: OCF twin vs. 16q + regression & 0.080 \\
Liquidity: dividend cut, twin added & 0.104 \\
Bankruptcy AUC: twin minus formula, all firms & 0.584 \\
\bottomrule
\end{tabular}

\end{table}
